\subsection{معرفی و انگیزه}
این ماژول در چارچوب پژوهش حاضر پیاده‌سازی و ارزیابی شده و فاقد مقاله مستقل جداگانه در ادبیات تحقیق است.

هرس نیمه‌ساختاریافته ماتریس‌های وزن ترنسفورمر، به‌ویژه در الگوهای ساختاری سخت‌افزاری \lr{2:4} و \lr{N:M}، مصالحه‌ای ایده‌آل میان کاهش حجم حافظه و حفظ سرعت پردازش در هسته‌های تانسوری مدرن فراهم می‌آورد. با این حال، سنجش اهمیت وزن‌ها بر اساس اندازه‌های قدرمطلقی صرف (\lr{Magnitude Pruning})، جهت‌گیری فضای ویژگی و برهم‌کنش همبستگی‌های ورودی را نادیده می‌گیرد. از سوی دیگر، محاسبه ماتریس کامل اطلاعات فیشر (\lr{Fisher Information Matrix}) مستلزم ضرب‌های ماتریسی با مرتبه $\mathcal{O}(d^4)$ است که برای لایه‌های ترنسفورمرهای امروزی عملاً غیرقابل اجراست.

برای حل این چالش، چارچوب فشرده‌سازی \lr{AST: Adaptive Semi-Structured Pruning with Empirical Fisher Calibration and PCA Coordinate Rotation} طراحی شده است. این روش ابتدا فعال‌سازی‌های ورودی را از طریق تحلیل مولفه‌های اصلی (\lr{PCA}) به یک دستگاه مختصات متعامد بدون همبستگی انتقال می‌دهد، سپس حساسیت ژاکوبی تجربی وزنی را در پایه‌های چرخش‌یافته محاسبه نموده و آن را به فضای مبنای اولیه تصویر می‌کند تا ماسک تنک ساختاریافته \lr{2:4} بدون برهم زدن تانسورهای اصلی استخراج گردد.
